% !TEX TS-program = xelatex
% !TEX encoding = UTF-8 Unicode

\documentclass{resume}
\usepackage{linespacing_fix}
\usepackage[hidelinks,pageanchor=false]{hyperref}

\setlist[itemize]{itemsep=0.15em}

\hypersetup{
  pdftitle={Zhongtao Rao - Machine Learning Engineer},
  pdfauthor={Zhongtao Rao}
}

\renewcommand{\contactInfo}[4]{%
  \centerline{\sffamily\normalsize
    #1 \quad\textperiodcentered\quad #2%
    \ifthenelse{\isempty{#3}}{}{\quad\textperiodcentered\quad #3}%
    \ifthenelse{\isempty{#4}}{}{\quad\textperiodcentered\quad #4}%
  }
  \vspace{0.7ex}
}

\begin{document}
\pagenumbering{gobble}

\name{Zhongtao Rao}
\centerline{\large Machine Learning Engineer | Recommender Systems \& Multimodal ML}
\vspace{0.7ex}
\contactInfo{\href{mailto:raozhongtao@gmail.com}{raozhongtao@gmail.com}}{+86 137 2777 2117}{}{}

\section{Profile}
Machine learning engineer specialising in recommender systems and multimodal representation learning, with experience in large-scale ranking, retrieval, user retention modelling and production software engineering. First author of a SIGIR 2026 paper; MPhil candidate graduating in January 2027.

\section{Education}
\textbf{The Hong Kong University of Science and Technology (Guangzhou)}
\datedline{MPhil in Data Science and Analytics}{Aug 2024 -- Jan 2027 (expected)}
\vspace{0.4em}
\textbf{South China University of Technology}
\datedline{BSc in Information and Computing Science}{Sep 2019 -- Jun 2023}

\section{Technical Skills}
\textbf{Languages:} Python, C/C++, Java, Swift\\
\textbf{ML frameworks \& deployment:} PyTorch, PEFT, ONNX\\
\textbf{Specialisms:} Recommendation ranking and retrieval, sequential recommendation, multimodal representation learning, contrastive learning, LoRA, multi-task learning, online experimentation

\section{Professional Experience}
\datedline{\textbf{Microsoft}}{Jul 2026 -- Present}
\textit{Applied Scientist Intern, Microsoft AI}

\vspace{0.3em}
\textbf{News ranking: model scaling and feature interactions}
\begin{itemize}
  \item Improved offline click AUC by \textbf{0.14\% relative to the baseline} by scaling a news-ranking model from 64M to approximately 95M parameters; trained from scratch on approximately \textbf{4 billion examples} over seven days.
  \item Replaced sequence pooling and the existing Transformer representation pipeline with HyFormer, using full self-attention over user behaviour sequences and a global token that cross-attends to the sequence.
  \item Designed a joint tokenizer to map sparse and dense features into eight non-sequential tokens; combined RankMixer feature interactions with MaskNet feature weighting and multi-task prediction heads.
\end{itemize}

\textbf{User retention prediction with multiple instance learning}
\begin{itemize}
  \item Developed a request-level multiple instance learning (MIL) model to learn article-level retention signals without assigning the same noisy request-level label to every article. Jointly predicted successful-session frequency over the next 24 hours and an exponentially decayed time-to-return target.
  \item Used user and context features to attend over article logits within each request; added a separate user tower to account for baseline return propensity and reduce the effect of user activity differences on learned content value.
  \item Exported an ONNX model for independent scoring of candidate articles and integrated retention predictions into ranking. Achieved offline AUCs of \textbf{0.8387} and \textbf{0.8064} for the frequency and time-to-return targets, respectively.
  \item Online experiments increased \textbf{successful sessions per user by 1.28\%}, time spent per user by \textbf{1.24\%}, and monetised content views per user by \textbf{1.06\%}.
\end{itemize}

\vspace{0.4em}
\datedline{\textbf{ByteDance}}{Apr 2026 -- Jul 2026}
\textit{Machine Learning Engineer Intern, Douyin Live}
\begin{itemize}
  \item Developed a traffic allocation strategy for mid- and long-tail revenue-generating livestream creators, setting high-quality traffic targets and applying targeted boosts across retrieval, pre-ranking and ranking.
  \item Identified high-potential creators using observed livestream engagement. Used a proportional--integral--derivative (PID) controller to adjust boosts based on the gap between actual and target traffic, with low-latency serving counters to stop boosts when targets were met.
  \item In a two-sided online experiment, increased streaming days per creator by \textbf{0.3657\%} and streaming duration per creator by \textbf{0.3806\%}.
\end{itemize}

\newpage
\section{Professional Experience (continued)}
\datedline{\textbf{Baidu}}{Nov 2025 -- Mar 2026}
\textit{Machine Learning Engineer Intern, Recommendation Strategy}
\begin{itemize}
  \item Built and deployed vision-language model embeddings for cold-start item-to-item (I2I) retrieval, addressing sparse interaction data and the semantic mismatch between separate BERT text and ViT image encoders.
  \item Constructed positive item pairs from adjacent clicks in user sequences and filtered weakly related popular items using co-occurrence scores to reduce popularity-driven noise.
  \item Fully fine-tuned \textbf{Qwen3-VL-2B} with in-batch contrastive learning on titles, speech transcripts, article text and images. Applied Matryoshka representation learning to reduce embedding dimensionality from \textbf{2,048 to 100}, lowering vector storage and retrieval costs.
  \item Online experiments on the cold-start feed increased first-hop click-through rate by \textbf{4.97\%}, second-hop distribution volume by \textbf{2.63\%} and completion rate by \textbf{1.53\%}, while reducing the quick-swipe rate by \textbf{1.70\%}.
\end{itemize}

\vspace{0.4em}
\datedline{\textbf{ByteDance}}{Jul 2023 -- Jul 2024}
\textit{Software Engineer, Feishu Mail \textnormal{(full-time)}}
\begin{itemize}
  \item Developed and maintained the email-reading module across Android and iOS, delivering product features and improving the user experience.
\end{itemize}

\section{Research Experience}
\textbf{Vision-Language Model Adaptation for Multimodal Recommendation}\\
\textit{First-author research, SIGIR 2026}
\begin{itemize}
  \item Adapted large vision-language models as unified item encoders to address misaligned image and text representations in multimodal recommendation. Encoded item images and structured text with hidden-state mean pooling.
  \item Designed item-level image--text contrastive learning and structure-aware soft labels derived from in-batch item similarities to preserve local semantic relationships.
  \item Proposed \textbf{Modality-Aware Mixture-of-Experts LoRA}, using modality-specific experts and routing to mitigate gradient conflicts during adaptation; used the resulting representations to initialise item ID embeddings for sequential recommendation.
  \item Evaluated Qwen2.5-VL-7B and Llama-3.2-11B-Vision on Amazon Beauty, Sports and Toys, with backbone comparisons and ablation studies. Improved \textbf{Hit@10 by 8.3\%} and \textbf{NDCG@10 by 10.4\%} on average; gains on long-tail items reached \textbf{74\%} and \textbf{95\%}, respectively.
\end{itemize}

\section{Publications \& Preprints}
\begin{itemize}
  \item \textbf{Z. Rao}, P. Zhou, D. Chong, Z. Chen, S. Wang, N. Tang.
  \textit{Structural and Disentangled Adaptation of Large Vision Language Models for Multimodal Recommendation.}
  \textbf{SIGIR 2026}. First author.

  \item Z. Chen, \textbf{Z. Rao}, W. Xiu, Y. Zhang, S. Cui, J. Chen, T. Zhao.
  \textit{LatentRemedy: Efficient Remedy For Non-Prefix KV Cache Reuse in Multi-Agent Systems.}
  \href{https://arxiv.org/abs/2610.18014}{arXiv:2610.18014}, 2026. Preprint; submitted to ACL 2027.

  \item Z. Chen, T. Wang, \textbf{Z. Rao}, H. Zhu, D. Cao, T. Zhao.
  \textit{SAGE: Structured Strategic Reasoning for Efficient LLM Game Playing.}
  \href{https://arxiv.org/abs/2609.34342}{arXiv:2609.34342}, 2026. Preprint; submitted to ICLR 2027.

  \item P. Zhou et al.\ (including \textbf{Z. Rao}).
  \textit{When Large Vision Language Models Meet Multimodal Sequential Recommendation: An Empirical Study.}
  \textbf{The Web Conference (WWW) 2025}. Oral presentation.
\end{itemize}

\end{document}
